\chapter{Modelo e implementación}\label{cap3:modelo}

\section{Formulación del problema}

Se consideran dominios adimensionales cuyas coordenadas están normalizadas por
la longitud de onda del láser ($\lambda = 638\,\mathrm{nm}$, láser de diodo
rojo). Bajo esta normalización, el número de onda
$k = 2\pi/\lambda$ definido en el Capítulo~\ref{cap2:marcos} toma el valor
adimensional $k = 2\pi$; salvo que se indique lo contrario, $k$ se usa en su
forma adimensional en el resto de este capítulo.

Para NB01 y NB02 se conserva el dominio adimensional $[0,1]$ y $[0,1]^2$,
respectivamente, porque esos experimentos validan la arquitectura contra
soluciones analíticas. NB03 requiere un dominio de propagación más amplio para
representar varios periodos ópticos y separar la pantalla de fase de la región
de observación. Por ello se define de forma independiente
$\Omega_{03}=[-10,10]\times[0,20]$ en coordenadas $(x/\lambda,z/\lambda)$,
manteniendo $k=2\pi$.

Se plantean tres experimentos sucesivos de complejidad creciente:

\subsection{Experimento 1 (NB01): Validación 1D}

La ecuación de Helmholtz unidimensional en $x \in [0,1]$ con condiciones de
Dirichlet:
%
\begin{equation}
    \frac{d^2 E}{dx^2} + k^2 E = 0, \quad
    E(x_j) = \cos(k x_j), \quad x_j \in \{0,\ 0.25,\ 0.5,\ 0.75,\ 1\},
    \label{eq:helmholtz1d}
\end{equation}
%
admite la solución analítica $E_\text{exacta}(x) = \cos(kx)$, utilizada para
validar la arquitectura y la función de pérdida. Se usan cinco puntos
conocidos en la pérdida de datos: los dos extremos $x=0$ y $x=1$ y tres puntos
interiores de supervisión. La
condición simétrica $E(0) = E(1) = 1$ colapsa la red hacia la solución
constante trivial durante el entrenamiento, ya que ambos extremos son
consistentes con $E \equiv 1$; añadir tres puntos intermedios conocidos de la
solución analítica rompe esa simetría y estabiliza la convergencia.

\subsection{Experimento 2 (NB02): Validación 2D con campo complejo}

La ecuación de Helmholtz 2D en $\Omega = [0,1]^2$ con condición de onda plana
diagonal en los cuatro bordes:
%
\begin{equation}
    E_\text{exacta}(x,y) = e^{i(k_x x + k_y y)}, \quad
    k_x = k_y = \frac{k}{\sqrt{2}},
    \label{eq:solucion2d}
\end{equation}
%
que satisface $k_x^2 + k_y^2 = k^2$. La condición de frontera se impone
separando la solución en sus partes real e imaginaria:
%
\begin{equation}
    E_\text{real}(x,y)\big|_{\partial\Omega} = \cos(k_x x + k_y y), \qquad
    E_\text{imag}(x,y)\big|_{\partial\Omega} = \sin(k_x x + k_y y).
\end{equation}

\subsection{Experimento 3 (NB03): Simulación de speckle óptico}
\label{cap3:nb03}

La ecuación de Helmholtz 2D en $\Omega_{03}$ con una pantalla de fase
aleatoria en la entrada $z=0$. La fase no se genera de forma independiente en
cada punto: se obtiene como un proceso gaussiano periódico con longitud de
correlación $\ell_\phi=0.5\lambda$ y desviación estándar $2$ rad. Esta
correlación evita que la condición de entrada quede dominada por componentes
evanescentes no observables en la región de propagación:
%
\begin{equation}
    E(x, 0) = \mathcal{P}_{|k_x|\leq k}
    \left\{e^{i\psi(x)}\right\}, \quad
    \psi(x)=\mathcal{G}_{\ell_\phi}[\xi(x)],
    \quad \xi(x)\sim\mathcal{N}(0,1),
    \label{eq:frontera_rugosa}
\end{equation}
%
donde $\mathcal{P}_{|k_x|\leq k}$ conserva los modos propagantes del espectro
angular. La derivada de entrada se fija mediante la condición saliente:
%
\begin{equation}
    \partial_z E(x,0)=
    \mathcal{F}^{-1}\!\left[
    i\sqrt{k^2-k_x^2}\,
    \widehat{E}(k_x,0)\right].
    \label{eq:derivada_saliente}
\end{equation}
%
En $x$ se imponen periodicidad de valor y de derivada. El dominio longitudinal
se divide en bloques cortos; cada bloque recibe el campo y la derivada de la
interfaz anterior, satisface Helmholtz en su interior y transfiere esas dos
cantidades al bloque siguiente. Esta formulación conserva la ecuación de
Helmholtz, pero evita tratar las fronteras restantes como bordes libres sin
condición física.

% ──────────────────────────────────────────────────────────────────────────────
\section{Arquitectura SIREN}

Se emplea una arquitectura SIREN con los parámetros consolidados a partir de la
validación 1D (NB01) y escalados para el caso 2D (NB02):

\begin{itemize}
    \item \textbf{Entrada:} coordenadas espaciales; $(x,y)\in[0,1]^2$ en NB02
    y $(x,z)\in\Omega_{03}$ en NB03.
    \item \textbf{Capas ocultas:} $L = 5$ capas de dimensión $d$ neuronas cada
    una ($d = 64$ en NB01, $d = 128$ en NB02/NB03).
    \item \textbf{Activación:} $\phi(z) = \sin(\omega_0 z)$. NB01/NB02 usan
    $\omega_0=1.0$; NB03 normaliza localmente las coordenadas de cada bloque y
    calibra $\omega_0$ como parte del experimento de alta frecuencia.
    \item \textbf{Salida:} capa lineal de dos neuronas $(E_\text{real}, E_\text{imag})$
    para el caso 2D, o una neurona $E$ para el caso 1D.
    \item \textbf{Parámetros totales:} $16{,}833$ (NB01) y $66{,}690$ (NB02).
\end{itemize}

La elección $\omega_0 = 1.0$ en NB01/NB02 se conserva porque esos casos ya
fueron validados en $[0,1]^2$. En NB03 el ancho físico del dominio es veinte
longitudes de onda y la red se entrena por bloques; por ello las coordenadas de
cada bloque se transforman a $[-1,1]^2$ y $\omega_0$ se trata como un
hiperparámetro calibrable. La selección definitiva se hará con error contra la
referencia de espectro angular, no únicamente con la pérdida de entrenamiento.

La inicialización sigue el esquema de \textcite{sitzmann2020siren}:
$\mathbf{W}_0 \sim \mathcal{U}(-1/n_\text{in}, 1/n_\text{in})$ para la primera
capa y $\mathbf{W}_\ell \sim \mathcal{U}(-\sqrt{6/d}/\omega_0,
\sqrt{6/d}/\omega_0)$ para capas intermedias. El objetivo es controlar la
escala de las activaciones y sus derivadas; no se supone que las activaciones
sean uniformes en $[-1,1]$.

% ──────────────────────────────────────────────────────────────────────────────
\section{Función de pérdida y calibración de pesos}

La función de pérdida total para el caso 2D es:
%
\begin{equation}
    \mathcal{L} = \mathcal{L}_\text{datos}
    + \lambda_\text{fís}\!\left(
    \mathcal{L}_{R_\text{real}} + \mathcal{L}_{R_\text{imag}}
    \right),
    \label{eq:loss_2d}
\end{equation}
%
con peso de física $\lambda_\text{fís} = 0.1$. Este valor fue calibrado
empíricamente mediante un experimento de ablación (véase la
Sección~\ref{sec:ablacion}). En esta implementación, la pérdida física suma
los residuos real e imaginario; con $\lambda = 1.0$ la corrida no completó el
entrenamiento, mientras que $\lambda = 0.01$ produjo un error ligeramente
menor pero corresponde a una restricción física de menor peso.

Los residuos de Helmholtz se calculan mediante diferenciación automática de
segundo orden \parencite{baydin2018autodiff}:
%
\begin{equation}
    R_\text{real}(\mathbf{r}) = \nabla^2 E_\text{real}(\mathbf{r})
    + k^2 E_\text{real}(\mathbf{r}), \qquad
    R_\text{imag}(\mathbf{r}) = \nabla^2 E_\text{imag}(\mathbf{r})
    + k^2 E_\text{imag}(\mathbf{r}),
\end{equation}
%
donde el Laplaciano se obtiene con \texttt{torch.autograd.grad} activando la
opción \texttt{create\_graph=True} para permitir diferenciación de segundo orden.
Es crítico que $k^2$ se convierta a tipo \texttt{float} de Python antes de
multiplicar por el tensor de la red, para evitar conflictos de tipo entre
\texttt{numpy.float64} y \texttt{torch.float32}.

Para NB03 se agrega la pérdida de la derivada longitudinal en la interfaz de
entrada, la periodicidad lateral y la transferencia entre bloques:
%
\begin{equation}
    \mathcal{L}_{03} =
    \mathcal{L}_{E,0}
    + \lambda_z\mathcal{L}_{\partial_z E,0}
    + \lambda_p\mathcal{L}_{\mathrm{per}}
    + \lambda_f\frac{\mathcal{L}_{R_\mathrm{real}}
    +\mathcal{L}_{R_\mathrm{imag}}}{k^4}.
    \label{eq:loss_nb03}
\end{equation}
%
La división por $k^4$ normaliza la escala del residuo de segundo orden. La
condición de derivada saliente no se usa como etiqueta interior: en el primer
bloque se calcula desde la pantalla radiativa y, en los bloques posteriores, se
transfiere desde la salida del bloque anterior. Así, la referencia angular se
reserva para evaluación independiente.

% ──────────────────────────────────────────────────────────────────────────────
\section{Estrategia de muestreo}

\subsection{Puntos interiores (colocación)}

Para NB02, los $N_c = 3{,}000$ puntos de colocación interiores se generan
mediante LHS \parencite{mckay1979lhs} usando
\texttt{scipy.stats.qmc.LatinHypercube} con dimensión $n = 2$. En NB03 los
puntos se muestrean dentro de cada bloque longitudinal, conservando cobertura
en $x$ y $z$ sin exigir una malla cartesiana global. La evaluación final se
realiza sobre la misma malla de la referencia angular.

\subsection{Puntos de frontera}

En NB01, la condición de Dirichlet se impone en los cinco puntos conocidos
descritos en la Sección~\ref{cap3:modelo} ($x = 0,\ 0.25,\ 0.5,\ 0.75,\ 1$).
En NB02, los $N_b = 300$ puntos de frontera por borde se distribuyen
uniformemente sobre cada uno de los cuatro lados del dominio $\partial\Omega$
(1{,}200 puntos en total). Para cada lado, se evalúa la solución exacta y se
construye el par $(\mathbf{r}^b_j,\, E(\mathbf{r}^b_j))$ utilizado en
$\mathcal{L}_\text{datos}$. NB03 sigue el esquema por bloques descrito en la
Sección~\ref{cap3:nb03}: la interfaz inferior recibe el campo y la derivada
saliente; las interfaces laterales imponen periodicidad de valor y derivada, y
las interfaces internas reciben únicamente la transferencia producida por la
PINN anterior. No se dejan bordes libres sin condición física.

% ──────────────────────────────────────────────────────────────────────────────
\section{Estrategia de optimización}

El entrenamiento sigue la estrategia bifásica descrita en el
Capítulo~\ref{cap2:marcos}, con los hiperparámetros consolidados en la
Tabla~\ref{tab:hiperparametros}.

\begin{table}[htbp]
\centering
\begin{threeparttable}
\caption{Hiperparámetros de entrenamiento consolidados.}
\label{tab:hiperparametros}
\begin{tabular}{lll}
\toprule
Hiperparámetro & Valor & Aplicación \\
\midrule
Tasa de aprendizaje Adam ($\eta$) & $10^{-3}$ & NB01, NB02 \\
Épocas máximas Adam             & 15{,}000  & NB01, NB02 \\
Paciencia (\textit{early stopping}) & 800 épocas & NB02 \\
$\delta$ mejora mínima          & $10^{-7}$ & NB02 \\
Iteraciones máximas L-BFGS      & 500 (NB01), 1{,}000 (NB02) & - \\
Historial L-BFGS                & 50 (NB01), 100 (NB02) & - \\
Semilla global (SEED)           & 42        & NB01, NB02 \\
Peso de física $\lambda_\text{fís}$ & 1.0 (NB01), 0.1 (NB02) & - \\
\bottomrule
\end{tabular}
\begin{tablenotes}\small
    \item Todos los experimentos se ejecutaron en GPU NVIDIA RTX 5050.
    La semilla favorece la reproducibilidad bajo las mismas condiciones de
    ejecución. NB01 usa $\lambda_\text{fís} = 1.0$
    porque un único residuo no duplica el peso de la física, a diferencia del
    caso 2D (real + imaginario); ver Sección~\ref{sec:ablacion}.
\end{tablenotes}
\end{threeparttable}
\end{table}

% ──────────────────────────────────────────────────────────────────────────────
\section{Métricas de validación}

\subsection{Error $L^2$ relativo (NB01 y NB02)}

El error $L^2$ relativo mide la discrepancia global entre la predicción de la
red $\hat{E}$ y la solución exacta $E_\text{exacta}$. Se evalúa sobre una malla
uniforme de $1{,}000$ puntos en $[0,1]$ para NB01 y de $100 \times 100$ puntos
en $[0,1]^2$ para NB02. Para NB01, que tiene salida escalar, se utiliza
%
\begin{equation}
    \varepsilon_{L^2} = \frac{\|\hat{E} - E_\text{exacta}\|_2}
    {\|E_\text{exacta}\|_2} \times 100\%.
    \label{eq:l2}
\end{equation}

En NB02, el código calcula esta cantidad por separado para las componentes real
e imaginaria y reporta su promedio aritmético:
%
\begin{equation}
\begin{aligned}
    \varepsilon_{L^2}^{(c)} & =
    \frac{\|\hat{E}_c - E_{c,\text{exacta}}\|_2}
    {\|E_{c,\text{exacta}}\|_2} \times 100\%,
    \quad c \in \{\mathrm{real},\mathrm{imag}\},\\
    \bar{\varepsilon}_{L^2} & =
    \frac{\varepsilon_{L^2}^{(\mathrm{real})}+
    \varepsilon_{L^2}^{(\mathrm{imag})}}{2}.
\end{aligned}
\label{eq:l2_componentes}
\end{equation}
%
Por tanto, $0.171\%$ es un promedio componente a componente y no la norma
relativa de un vector complejo concatenado.

Para esta validación no se adopta un criterio universal de tolerancia de error;
se fijó
$\varepsilon_{L^2} < 5\%$ como umbral ingenieril propio de la tesis. Este
valor resulta, además, considerablemente menor que el orden de magnitud
reportado en implementaciones comparables de PINNs para la ecuación de
Helmholtz identificadas posteriormente (Tabla~\ref{tab:comparativa}).

% ──────────────────────────────────────────────────────────────────────────────
\section{Implementación modular}

El código fuente está organizado en módulos reutilizables en \texttt{src/}:

\begin{description}
    \item[\texttt{models.py}] Clase \texttt{PINN\_2D\_SIREN} que implementa
    la arquitectura SIREN con la inicialización de \textcite{sitzmann2020siren}.
    El método \texttt{forward(x)} aplica la activación sinusoidal capa por capa
    y retorna el campo complejo como un tensor de dos columnas
    $(E_\text{real}, E_\text{imag})$.

    \item[\texttt{losses.py}] Función \texttt{pinn\_loss\_2d(model, xy\_int,
    xy\_bc, E\_bc, k, lam)} que calcula la pérdida total~\eqref{eq:loss_2d}.
    El cálculo del Laplaciano utiliza \texttt{torch.autograd.grad} con
    \texttt{create\_graph=True} para habilitar la diferenciación de segundo orden.
    El escalar $k$ se convierte explícitamente a \texttt{float(k)} antes de
    multiplicar por el tensor de la red, evitando conflictos de dtype entre
    \texttt{numpy.float64} y \texttt{torch.float32}.

    \item[\texttt{training.py}] Loop de entrenamiento que integra las fases
    Adam y L-BFGS con \textit{early stopping}.

    \item[\texttt{utils.py}] Métricas ($L^2$, $R^2$, RMSE, MAE), funciones de
    muestreo LHS y rutinas de visualización.
\end{description}

Los experimentos están reproducidos en los notebooks
\texttt{01\_pinn\_helmholtz\_1d\_validation.ipynb},
\texttt{02\_pinn\_helmholtz\_2d\_complex\_field.ipynb} y
\texttt{03\_pinn\_optical\_speckle\_simulation.ipynb}.
Los resultados numéricos están archivados en formato JSON en \texttt{results/}.
